\section*{Chapter \ref{ch:in:multi-manager-platforms} --- Multi-Manager Platforms}

\begin{solution}{exo:in:multi-manager-platforms:1}
Employees $6\,140/1\,625=3.78$ times; RAUM $571.1/181.5=3.15$ times.
\end{solution}

\begin{solution}{exo:in:multi-manager-platforms:2}
$\ln2/0.2=3.5$ years.
\end{solution}

\begin{solution}{exo:in:multi-manager-platforms:3}
$360\times5=1\,800$ people in teams, about 26\% of 7\,000: roughly three quarters of the staff run the common machinery.
\end{solution}

\begin{solution}{exo:in:multi-manager-platforms:4}
$183/287=64\%$ in 2016 and $869/2\,063=42\%$ in 2026. Most new positions were outside advisory functions: technology,
data, risk and operations grew faster than the teams.
\end{solution}

\begin{solution}{exo:in:multi-manager-platforms:5}
53.5\% after three years; half are closed after 3.6 years.
\end{solution}

\begin{solution}{exo:in:multi-manager-platforms:6}
At 1.25 years: 94.5\% of the loose ladder's teams are running against 54.2\% of the tight ladder's.
\end{solution}

\begin{solution}{exo:in:multi-manager-platforms:7}
The median falls from 3.6 years to about 0.9 years, and 98.6\% are closed within ten years: with the same Sharpe ratio,
a higher volatility makes the same loss levels a smaller number of standard deviations, so they are reached sooner.
The ladder is set in capital, not in units of the team's risk.
\end{solution}

\begin{solution}{exo:in:multi-manager-platforms:8}
The platform's alternative is to keep bad teams longer: under the tight ladder most closed teams (61\%) were not
skilled. The share of skilled teams among those closed is the cost of speed; the right ladder weighs it against the losses of
keeping bad teams, and its level should scale with each team's volatility.
\end{solution}

\begin{solution}{pb:in:multi-manager-platforms:1}
\begin{enumerate}
\item Team: strategy, people, trading decisions within limits. Platform: capital, data, execution, technology, risk,
compliance, operations, financing.
\item 1\,625 and 6\,140; 3.78 times.
\item RAUM per employee fell by 17\% (a factor of 0.83); advisory share from 63\% to 47\%.
\item More than 360 investment teams; 165 teams in six strategies.
\item No primary source: press counts are not filed or published by the platforms.
\item Time from start to closure; closures in a year over the average number of teams; median $\ln2/\lambda$.
\item Tight: halve at 5\%, close at 7.5\%; loose: 10\% and 15\%; volatility 10\%.
\item 1.5 years and 3.6 years.
\item Tight: 94.4\% and 74.8\%; loose: 41.4\% and 15.3\%.
\item It has no gains yet, so its peak is its starting capital and any early loss is a drawdown.
\item 20\% and 5\%.
\item 35\% and 3\%.
\item 39\% and 14\%.
\item It lengthens it: without halving the median is 0.7 years instead of 1.5.
\item It stays within 10\% of 1.5 years (1.6 years with four times as many steps).
\item Tight ladder: 1.5 and 3.6 years; skilled share of closed teams 39\% (tight) and 14\% (loose).
\item The ladder's levels, whether they scale with the strategy's volatility, and how capital is added after gains.
\item Analysts are employed by the team, and their jobs usually end with it.
\item Reviews on other criteria, capital added after gains, teams leaving on their own.
\item With a Sharpe ratio of 0.5 to 1.0 under a tight ladder she should expect two to four years; she should negotiate
the ladder or the size so that her volatility makes the stop levels several standard deviations away.
\end{enumerate}
\end{solution}

\begin{solution}{iq:in:multi-manager-platforms:1}
A fund whose capital is run by many independent teams, each with its own limits and pay, on common infrastructure.
In a team you run a small business under the platform's limits and are paid on your team's P\&L; at a single-strategy
fund you contribute to one book and are paid on the firm's result.
\iqlookfor{the team as the unit of pay and risk.}
\end{solution}

\begin{solution}{iq:in:multi-manager-platforms:2}
The year's P\&L is about normal with mean 10\% and standard deviation 10\%: $P(X\le-5\%)=\Phi(-1.5)=6.7\%$.
\iqlookfor{mean, standard deviation, a z-score.}
\end{solution}

\begin{solution}{iq:in:multi-manager-platforms:3}
A first drawdown is weak evidence: halving limits the loss if the team is bad while keeping it alive if it was
unlucky, and doubles the distance to the stop in units of the remaining risk.
\iqlookfor{evidence against luck, and the value of a smaller position.}
\end{solution}

\begin{solution}{iq:in:multi-manager-platforms:4}
Survivors are selected on realised P\&L, so they are better than the teams hired, but mixed with lucky ones; the
surviving pool's average skill rises with time, and its measured returns overstate its skill (survivorship bias,
Book~7, chapter~3).
\iqlookfor{selection on outcomes, and its bias.}
\end{solution}

\begin{solution}{iq:in:multi-manager-platforms:5}
Small at first: until gains build room below the peak, a normal loss can reach the stop. Scale up as the cushion
grows, keeping the stop several standard deviations of expected P\&L away.
\iqlookfor{sizing to distance from the stop.}
\end{solution}

\begin{solution}{iq:in:multi-manager-platforms:6}
From the start: the probability that $\mu t+\sigma W_t$ ever reaches $-d$ is $e^{-2\mu d/\sigma^2}$ (Book~4). From the peak:
the drawdown is a reflected Brownian motion with drift towards zero, which is recurrent and so reaches any level $d$ with
probability one. A ladder measured from the peak closes every team eventually; only its timing depends on skill.
\iqlookfor{the two probabilities, and why the peak-based stop always triggers.}
\end{solution}
